\subsection{Sealed Bids Method (Cont.)}


\content{
\begin{example} Four small graphic design companies are
merging operations to become one larger corporation. In this
merger, there are a number of issues that need to be
settled. Each company is asked to place a monetary value (in
thousands of dollars) on each issue: \\

\begin{tabular}{|l|l|l|l|l|}
\hline
& Super Designs & DesignByMe & LayoutPros & Graphix \\
\hline
Company name & \$5 & \$3 & \$3 & \$6 \\
\hline
Company location & \$8 & \$9 & \$7 & \$6 \\
\hline
CEO & \$10 & \$5 & \$6 & \$7 \\
\hline
Chair of the Board & \$7 & \$6 & \$6 & \$8 \\
\hline
\end{tabular}
\end{example}
}{% presentation
\vspace{3in}
}{% nopresentation
\vspace{2in}
}{% comments
}

\content{

        \begin{example} Fair division does not always have to be used for 
        items of value. It can also be used to divide undesirable items. 
        Suppose Chelsea and Mariah are sharing an apartment, and need to 
        split the chores for the household. They list the chores, assigning 
        a negative dollar value to each item; in other words, the amount 
        they would pay for someone else to do the chore (a per week amount).
        We will assume, however, that they are committed to doing all 
        the chores themselves and not hiring a maid. Find a fair division if the
        bids are given as shown below. \\

        \begin{tabular}{|l|l|l|}
        \hline
        & Chelsea & Mariah \\
        \hline
        Vacuuming & -\$10 & -\$8 \\
        \hline
        Cleaning bathroom  & -\$14 & -\$20 \\
        \hline
        Doing Dishes & -\$4 & -\$6 \\
        \hline
        Dusting & -\$6 & -\$4 \\
        \hline
        \end{tabular}
        \end{example} 
}{% presentation
\vspace{3in}
}{% nopresentation
\vspace{2in}
}{% comments
}

\content{
\begin{example} Four heirs (A, B, C, and D) must fairly
divide an estate consisting of two items - a desk and a
vanity - using the method of sealed bids. The players' bids
(in dollars) are: \\

\begin{tabular}{|l|l|l|l|l|}
\hline
& A & B & C & D \\
\hline
Desk & 320 & 240 & 300 & 260 \\
\hline
Vanity & 220 & 140 & 200 & 180 \\
\hline
\end{tabular}
\end{example}
}{% presentation
\vspace{3in}
}{% nopresentation
\vspace{2in}
}{% comments
}

\content{
\begin{example} After living together for a year, Emily,
Kayla, and Kendra have decided to go their separate ways.
They have several items they bought together to divide, as
well as some moving-out chores. The values each place are
shown below. Find the final allocation. \\

\begin{tabular}{|l|l|l|l|}
\hline
& Emily & Kayla & Kendra \\
\hline
Dishes & 20 & 30 & 40 \\
\hline
Vacuum cleaner & 100 & 120 & 80 \\
\hline
Dining table & 100 & 80 & 130 \\
\hline
Detail cleaning & -70 & -40 & -50 \\
\hline
Cleaning carpets & -30 & -60 & -50 \\
\hline
\end{tabular}
\end{example}
}{% presentation
\vspace{3in}
}{% nopresentation
\vspace{3in}
}{% comments
}
